% GENERATED FILE. Edit canonical lessons, then run npm run generate:books.
% canonical-source-hash: fnv1a64:ec6ee70fa659743e
% canonical-lessons: ZH-C230-shenshang, ZH-W230-shen-body, ZH-C230-tizhong, ZH-W230-ti-body, ZH-C230-duzi, ZH-C230-bizi, ZH-C230-tui

\chapter{Head to Toe}
\label{ch:zh-head-to-toe}
\hlchaptermodality{230}

\begin{chapteropening}
\textbf{By the end of this chapter:} \emph{I can say what I have on me and what I weigh, and talk about my tummy, my nose and my legs.}

Complete the last lesson of chapter 230: I can say what I have on me and what I weigh, and talk about my tummy, my nose and my legs.
\end{chapteropening}

\section[\texorpdfstring{\zh{身上} (shēnshang)}{身上 (shēnshang)}]{\zh{身上} (shēnshang) — on your person; on one's body}

\label{lesson:ZH-C230-shenshang}



\begin{quote}
Before the new one: say the Mandarin for to scratch, then the Mandarin for to stroke.
\end{quote}

\subsection*{You'll want to know: \zh{身上}}
\textbf{\zh{身上}} — \emph{shēnshang} — \textquotedblleft{}on your person; on one's body\textquotedblright{}.

\begin{quote}
\textbf{\zh{我身上没有钱}} \emph{wǒ shēnshang méiyǒu qián} — I haven't got any money on me.
\end{quote}

The new character, in pieces: \zh{身} — seven strokes: a short fall to the left, a short vertical, a bar that turns down into a long hook, two short bars inside, a long stroke rising to the right, and a long stroke falling to the left. The next lesson puts it in your hand.

\begin{cousinweb}
\textbf{\zh{身上}} (\emph{shēnshang}) is built from \zh{身} 'body' and \zh{上} 'on', both inherited from Middle Chinese, \emph{syin} and \emph{dzyangH} (Guangyun, in Baxter's transcription), which glosses the first 'oneself' and 'the body'; in the word the second loses its tone.
\end{cousinweb}

\begin{grammarlens}[title={What you've built}]
One more word for describing things.
\end{grammarlens}

\subsection*{Your turn}
\begin{itemize}
\raggedright
  \item \emph{Say it:} \emph{shēnshang}
  \item \emph{Say it:} \emph{shēnshang}, once more
  \item \emph{Say it:} \emph{wǒ shēnshang méiyǒu qián} — I haven't got any money on me.
  \item \emph{Recall:} say the Mandarin for to scratch, then the Mandarin for to stroke, then say \emph{shēnshang} again
\end{itemize}

\begin{tcolorbox}[breakable,colback=teal!4,colframe=teal!35!black,title={Before you move on}]
What does \zh{身上} mean? (\textquotedblleft{}On your person; on one's body\textquotedblright{}.) Say it once more.
\end{tcolorbox}
\section[\texorpdfstring{\zh{身} (shēn)}{身 (shēn)}]{\zh{身} — a frame with two bars, crossed by two long strokes}

\label{lesson:ZH-W230-shen-body}



\begin{quote}
The word first, then the character inside it.

\begin{itemize}
\raggedright
  \item \emph{Recall:} say \emph{shēnshang} — \textbf{R1}, one lesson back, before it decays
\end{itemize}

\begin{quote}
\zh{上} \zh{月}
\end{quote}

\emph{Read it:} each one aloud, with nothing said first — \emph{shàng}, \emph{yuè}
\end{quote}

\subsection*{Script — a frame with two bars, crossed by two long strokes}
\hlblockfigure{\detokenize{figures/ZH-W230-shen-body-filmstrip.pdf}}{How \zh{身} is written, stroke by stroke}

\begin{quote}
\zh{身}
\end{quote}

\textbf{One piece, seven strokes.} Count them in three groups.

\textbf{The frame, three strokes.} A short stroke falling to the left at the top. A short vertical on the left. Then a bar from its top that runs right and turns down into a long stroke to the foot, hooking left.

\textbf{Inside, two short bars,} one under the other.

\textbf{Across it all, two long strokes.} A stroke rising from the far left to the right, through the frame. Then a long stroke from the upper right falling down to the lower left.

3 + 2 + 2 = seven.

\subsection*{Your turn}
\begin{itemize}
\raggedright
  \item \emph{Write it:} \zh{身}, all seven, with the model in front of you
  \item \emph{Cover:} write it again from memory, and count the pieces, not the strokes
\end{itemize}

\begin{tcolorbox}[breakable,colback=teal!4,colframe=teal!35!black,title={Before you move on}]
How many strokes in \zh{身}? (\textbf{Seven}.) How many short bars sit inside the frame? (\textbf{Two}.) What comes last? (\textbf{The long stroke falling to the left}.)

Source: \href{https://raw.githubusercontent.com/chanind/hanzi-writer-data/68d10a4b21150cae5e1ebbd223eed289cf32d90c/data/\%E8\%BA\%AB.json}{Hanzi Writer Data \zh{身}, pinned stroke paths and medians}.
\end{tcolorbox}
\section[\texorpdfstring{\zh{体重} (tǐzhòng)}{体重 (tǐzhòng)}]{\zh{体重} (tǐzhòng) — body weight}

\label{lesson:ZH-C230-tizhong}



\begin{quote}
Before the new one: say the Mandarin for to stroke, then the Mandarin for on your person.
\end{quote}

\subsection*{You'll want to know: \zh{体重}}
\textbf{\zh{体重}} — \emph{tǐzhòng} — \textquotedblleft{}body weight\textquotedblright{}.

\begin{quote}
\textbf{\zh{你的体重是多少}} \emph{nǐ de tǐzhòng shì duōshao} — How much do you weigh?
\end{quote}

The new character, in pieces: \zh{体} — seven strokes: the person piece \zh{亻}, then \zh{本}. The next lesson puts it in your hand.

\begin{cousinweb}
\textbf{\zh{体重}} (\emph{tǐzhòng}) is built from \zh{体} 'body' and \zh{重} 'heavy', both inherited from Middle Chinese, \emph{thejX} and \emph{drjowngX} (Guangyun, in Baxter's transcription), which glosses the first 'the body' and the second 'much' and 'thick'.
\end{cousinweb}

\begin{grammarlens}[title={What you've built}]
One more word for describing things.
\end{grammarlens}

\subsection*{Your turn}
\begin{itemize}
\raggedright
  \item \emph{Say it:} \emph{tǐzhòng}
  \item \emph{Say it:} \emph{tǐzhòng}, once more
  \item \emph{Say it:} \emph{nǐ de tǐzhòng shì duōshao} — How much do you weigh?
  \item \emph{Recall:} say the Mandarin for to stroke, then the Mandarin for on your person, then say \emph{tǐzhòng} again
\end{itemize}

\begin{tcolorbox}[breakable,colback=teal!4,colframe=teal!35!black,title={Before you move on}]
What does \zh{体重} mean? (\textquotedblleft{}Body weight\textquotedblright{}.) Say it once more.
\end{tcolorbox}
\section[\texorpdfstring{\zh{体} (tǐ)}{体 (tǐ)}]{\zh{体} — the person piece, then \zh{本}}

\label{lesson:ZH-W230-ti-body}



\begin{quote}
The word first, then the character inside it.

\begin{itemize}
\raggedright
  \item \emph{Recall:} say \emph{tǐzhòng} — \textbf{R1}, one lesson back, before it decays
\end{itemize}

\begin{quote}
\zh{本} \zh{你}
\end{quote}

\emph{Read it:} each one aloud, with nothing said first — \emph{běn}, \emph{nǐ}
\end{quote}

\subsection*{Script — the person piece, then \zh{本}}
\hlblockfigure{\detokenize{figures/ZH-W230-ti-body-filmstrip.pdf}}{How \zh{体} is written, stroke by stroke}

\begin{quote}
\zh{体}
\end{quote}

\textbf{Two pieces, left before right.}

\textbf{The person piece, \zh{亻}, two strokes,} as in \zh{你} and \zh{他}.

\textbf{Then \zh{本}, five strokes,} as in \zh{本子}: a bar across, a long vertical, a stroke falling to the left and one falling to the right, then a short bar low on the vertical.

2 + 5 = seven.

\subsection*{Your turn}
\begin{itemize}
\raggedright
  \item \emph{Write it:} \zh{体}, all seven, with the model in front of you
  \item \emph{Cover:} write it again from memory, and count the pieces, not the strokes
\end{itemize}

\begin{tcolorbox}[breakable,colback=teal!4,colframe=teal!35!black,title={Before you move on}]
How many strokes in \zh{体}? (\textbf{Seven}.) What are its two pieces? (\textbf{\zh{亻} and \zh{本}}.) What comes last? (\textbf{The short bar near the foot}.)

Source: \href{https://raw.githubusercontent.com/chanind/hanzi-writer-data/68d10a4b21150cae5e1ebbd223eed289cf32d90c/data/\%E4\%BD\%93.json}{Hanzi Writer Data \zh{体}, pinned stroke paths and medians}.
\end{tcolorbox}
\section[dùzi]{dùzi — belly; tummy}

\label{lesson:ZH-C230-duzi}



\begin{quote}
Before the new one: say the Mandarin for on your person, then the Mandarin for body weight.
\end{quote}

\subsection*{You'll want to know: dùzi}
\emph{dùzi} — \textquotedblleft{}belly; tummy\textquotedblright{}.

\begin{quote}
\emph{wǒ dùzi téng} — My tummy hurts.
\end{quote}

Its characters come later; for now it is a word for the ear.

\begin{cousinweb}
\emph{dùzi} is built from \emph{dù} 'belly' and the noun ending \zh{子}, both inherited from Middle Chinese, \emph{duX} and \emph{tsiX} (Guangyun, in Baxter's transcription), which glosses the first 'the belly'.
\end{cousinweb}

\begin{grammarlens}[title={What you've built}]
One more word for describing things.
\end{grammarlens}

\subsection*{Your turn}
\begin{itemize}
\raggedright
  \item \emph{Say it:} \emph{dùzi}
  \item \emph{Say it:} \emph{dùzi}, once more
  \item \emph{Say it:} \emph{wǒ dùzi téng} — My tummy hurts.
  \item \emph{Recall:} say the Mandarin for on your person, then the Mandarin for body weight, then say \emph{dùzi} again
\end{itemize}

\begin{tcolorbox}[breakable,colback=teal!4,colframe=teal!35!black,title={Before you move on}]
What does dùzi mean? (\textquotedblleft{}Belly; tummy\textquotedblright{}.) Say it once more.
\end{tcolorbox}
\section[bízi]{bízi — nose}

\label{lesson:ZH-C230-bizi}



\begin{quote}
Before the new one: say the Mandarin for body weight, then the Mandarin for a tummy.
\end{quote}

\subsection*{You'll want to know: bízi}
\emph{bízi} — \textquotedblleft{}nose\textquotedblright{}.

\begin{quote}
\emph{tā de bízi hěn gāo} — He has a high-bridged nose.
\end{quote}

Its characters come later; for now it is a word for the ear.

\begin{cousinweb}
\emph{bízi} is built from \emph{bí} 'nose' and the noun ending \zh{子}, both inherited from Middle Chinese, \emph{bjijH} and \emph{tsiX} (Guangyun, in Baxter's transcription), which explains the first, after an older book, as drawing in breath. It is listed with a departing tone, so the rising tone of bí is not the regular outcome.
\end{cousinweb}

\begin{grammarlens}[title={What you've built}]
One more word for describing things.
\end{grammarlens}

\subsection*{Your turn}
\begin{itemize}
\raggedright
  \item \emph{Say it:} \emph{bízi}
  \item \emph{Say it:} \emph{bízi}, once more
  \item \emph{Say it:} \emph{tā de bízi hěn gāo} — He has a high-bridged nose.
  \item \emph{Recall:} say the Mandarin for body weight, then the Mandarin for a tummy, then say \emph{bízi} again
\end{itemize}

\begin{tcolorbox}[breakable,colback=teal!4,colframe=teal!35!black,title={Before you move on}]
What does bízi mean? (\textquotedblleft{}Nose\textquotedblright{}.) Say it once more.
\end{tcolorbox}
\section[tuǐ]{tuǐ — leg}

\label{lesson:ZH-C230-tui}



\begin{quote}
Before the new one: say the Mandarin for a tummy, then the Mandarin for a nose.
\end{quote}

\subsection*{You'll want to know: tuǐ}
\emph{tuǐ} — \textquotedblleft{}leg\textquotedblright{}.

\begin{quote}
\emph{wǒ shēnshang téng tuǐ yě téng} — I ache all over, and my legs too.
\end{quote}

Its characters come later; for now it is a word for the ear.

\begin{cousinweb}
\emph{tuǐ} is inherited from Middle Chinese \emph{thwojX} (Guangyun, in Baxter's transcription), glossed 'the thigh'; the dictionary lists the character as the everyday form of another.
\end{cousinweb}

\begin{grammarlens}[title={What you've built}]
One more word for describing things.
\end{grammarlens}

\subsection*{Your turn}
\begin{itemize}
\raggedright
  \item \emph{Say it:} \emph{tuǐ}
  \item \emph{Say it:} \emph{tuǐ}, once more
  \item \emph{Say it:} \emph{wǒ shēnshang téng tuǐ yě téng} — I ache all over, and my legs too.
  \item \emph{Recall:} say the Mandarin for a tummy, then the Mandarin for a nose, then say \emph{tuǐ} again
\end{itemize}

\begin{tcolorbox}[breakable,colback=teal!4,colframe=teal!35!black,title={Before you move on}]
What does tuǐ mean? (\textquotedblleft{}Leg\textquotedblright{}.) Say it once more.
\end{tcolorbox}
